\section{Properness, club witnesses and parameter invariance}
\label{sec:p06-properness-parameter-invariance}

An iteration construction may produce masters for models containing its
full datum, while properness tests all models containing the bare poset.
The conversion is not downward inheritance of a master's genericity:
it is a small-trace absoluteness argument. This host proof uses the
master and height conventions of
\cref{def:p06-pfm-hchi,def:p06-pfm-master-condition,def:p06-pfm-proper}.
For the all-model characterization in a primary exposition, see
\cite[pp.~48--49]{Dow1988ElementarySubmodels}. The argument here supplies
the club/code and parameter-removal construction at the present hereditary
threshold, rather than reproducing that exposition's referred proof route.

\begin{hbconvention}{def:p06-ppi-standing}{The poset and exact hereditary threshold}
Work in \(\mathsf{ZFC}\), with nonempty set poset
\(\mathbf P=(P,\leq_P,1_P)\), smaller stronger. Put
\[
 \kappa=\max\{\aleph_1,|\operatorname{tc}(\{\mathbf P\})|\},
 \qquad X=\omega\cup P\cup\mathcal P(P).
\]
A book-admissible height is regular \(\chi>2^\kappa\).
The trace universe includes all dense subsets, not only their names or a
selected list. The natural-number part makes every model trace countably
infinite, including when the poset is finite. There is no inaccessible,
GCH, ccc, closure or generic-existence hypothesis in this bridge.
\end{hbconvention}

\begin{hbdefinition}{def:p06-ppi-good-trace}{The master test on a small trace}
Let \(\mathcal D_P\) be the set of all dense subsets of \(P\).
For \(Y\in[X]^\omega\), let \(\operatorname{Good}_P(Y)\) mean
\[
 (\forall p\in P\cap Y)(\exists q\in P)\bigl(q\leq_P p\ \land
  (\forall D\in\mathcal D_P\cap Y)(\forall r\in P)
   \bigl(r\leq_P q\ \Longrightarrow
    (\exists d\in D\cap Y)(\exists s\in P)\ s\leq_P r,d\bigr)\bigr).
\]
The same \(q\) works for every dense set in the trace; it need not
belong to \(Y\). Compatibility, not a model condition strengthening
\(r\), is requested. A club on \([Z]^\omega\) is unbounded under
inclusion and closed under increasing countable unions, as in
\cref{thm:p06-pfm-hulls}.
\end{hbdefinition}

\begin{hblemma}{thm:p06-ppi-small-absoluteness}{proved=annex:p06.ppi.small-absoluteness}{The trace and every witness code fit at every admissible height}
The sets \(X\), \([X]^\omega\), \([X]^{\leq\omega}\), and
\(\mathcal D_P\) belong to every book-admissible \(H_\chi\).
So does every individual club \(C\subseteq[X]^\omega\) and every
individual function \(g:[X]^{<\omega}\to[X]^{\leq\omega}\).
Their relevant hereditary sizes are at most \(2^\kappa\).
Countability, clubhood, closure under \(g\), and \(\operatorname{Good}_P\)
are absolute between \(H_\chi\) and the ground universe. This does not
assume that the collection of all such codes or clubs belongs to \(H_\chi\).
For countable \(N\prec H_\chi\) containing \(\mathbf P\), its trace
\(N\cap X\) is countably infinite, and \(\operatorname{Good}_P(N\cap X)\)
is exactly the book's master-production test for that \(N\).
\end{hblemma}

\begin{hblemma}{thm:p06-ppi-club-codes}{proved=annex:p06.ppi.club-codes}{Finite-input countable-output codes for clubs}
Let \(Z\) be infinite. For
\(g:[Z]^{<\omega}\to[Z]^{\leq\omega}\), the collection
\[
 C_g=\{Y\in[Z]^\omega:(\forall a\in[Y]^{<\omega})\ g(a)\subseteq Y\}
\]
is club. Conversely every club \(C\subseteq[Z]^\omega\) contains
\(C_g\) for some such \(g\). The converse selects monotone club
supersets of finite inputs, then uses an increasing enumeration of each
countable closure point; no unspecified club-coding theorem is imported.
\end{hblemma}

\begin{hblemma}{thm:p06-ppi-trace-projection}{proved=annex:p06.ppi.trace-projection}{A club of large good models supplies a small good-trace club}
Suppose at some regular \(\Theta>2^\kappa\), a club
\(C\subseteq[H_\Theta]^\omega\) consists of elementary models
containing \(\mathbf P\), each satisfying the master-production test.
Then \(\{Y\in[X]^\omega:\operatorname{Good}_P(Y)\}\) contains a club.
The proof constructs a club contained in the image of the trace map;
it does not assert that an arbitrary projected image is itself closed.
Large auxiliary parameters can occur in the original models while the
induced code on \(X\) remains hereditarily small.
\end{hblemma}

\begin{hblemma}{thm:p06-ppi-all-model-recognition}{proved=annex:p06.ppi.all-model-recognition}{Every bare-poset elementary model selects its own code}
If the good traces contain a club, then for every regular
\(\chi>2^\kappa\), every countable \(N\prec(H_\chi,\in)\) required
to contain \(\mathbf P\) but no auxiliary parameter, and every
\(p\in P\cap N\), a master
\(q\leq_P p\) exists. Neither the original large parameter nor the
original club witness is required to belong to \(N\).
\end{hblemma}

\begin{hbtheorem}{thm:p06-ppi-equivalence}{proved=annex:p06.ppi.equivalence}{Club-good-model and parameter-enriched tests equal the all-model test}
The following are equivalent:
\begin{enumerate}
\item \(\mathbf P\) is proper in \cref{def:p06-pfm-proper}, at every
      regular \(\chi>2^\kappa\) and every poset-containing model.
\item The good traces on \(X\) contain a club.
\item At some regular \(\Theta>2^\kappa\), there is a club
      of elementary \(N\prec H_\Theta\) containing \(\mathbf P\),
      all satisfying the master-production test.
\item At some regular \(\Theta>2^\kappa\) and for some
      fixed parameter \(A\in H_\Theta\), every countable elementary
      \(N\prec H_\Theta\) containing \(\mathbf P,A\) has the
      master-production property for every \(p\in P\cap N\).
\end{enumerate}
In the last item \(A\) may be the entire iteration datum. One such
large-height witness implies the all-model conclusion at \emph{every}
book-admissible height, even heights not containing \(A\).
\end{hbtheorem}

\begin{hbcorollary}{thm:p06-ppi-cs-proper}{proved=annex:p06.ppi.cs-proper}{Properness of every countable-support stage}
Under the hypotheses of \cref{thm:p06-pil-extension}, every initial
segment and the full countable-support iteration are proper in the
book's bare-poset/all-model formulation. The premise is master production
with the full iteration datum, \cref{thm:p06-pil-model-masters}; the
equivalence theorem removes the auxiliary datum from the model test.
\end{hbcorollary}

\begin{hbtheorem}{thm:p06-ppi-no-downward-inheritance}{proved=annex:p06.ppi.no-downward-inheritance}{A master for a larger model need not be a master for the smaller one}
There are a proper atomic poset, countable \(N\prec M\prec H_\chi\)
containing it, and a condition which is master for \(M\) but not for
\(N\). Thus the parameter-invariance proof cannot be replaced by taking
a larger parameter-containing model and inheriting its chosen master.
\end{hbtheorem}

\begin{proofboundary}{def:p06-ppi-source-boundary}{Parameters and preservation hypotheses}
The club/all-model parameter bridge holds in \(\mathsf{ZFC}\) at the
stated hereditary threshold. It compares master-production formulations;
it is not a PFA/MM consistency construction. The support-controlled
proper-iteration lemma remains the substantive preservation premise.
The original Abraham and Shelah presentations have their own proof
routes, not supplied here as unchanged reproductions.
\end{proofboundary}
